\chapter{專案簡介}

\section{專案目標}

本專案旨在實作一個基於 C11 標準的 SIC/XE 組合語言組譯器（Assembler）。
其核心功能是將符合 SIC/XE 指令集架構的組合語言原始碼（\texttt{.asm}）轉換為計算機可執行之目的程式（\texttt{.obj}）。
本實作不僅支援標準的 SIC 指令，亦支援 XE 架構的擴充功能，如 Format 4 指令、相對定址模式等。

\section{核心特色}

\begin{itemize}
    \item \textbf{完整指令集支援}：支援 Format 1、2、3 及 4（前綴 \lstinline|+|）之指令解析。
    \item \textbf{定址模式自動處理}：自動計算 PC 相對定址（PC-Relative）與 Base 相對定址（Base-Relative）之位
          移量（Displacement）。
    \item \textbf{嚴謹的語法解析}：採用有限狀態機（FSM）模型進行詞法分析，確保能精準捕捉 Label、Mnemonic 與 Operands，
          並具備基礎語法錯誤偵測。
    \item \textbf{跨平台建置}：使用 \texttt{CMake} 進行專案自動化建置，支援 Linux、macOS 與 Windows 環境。
\end{itemize}

\section{開發環境與工具}

\begin{itemize}
    \item \textbf{程式語言}：C（C11 標準）。
    \item \textbf{建置系統}：CMake 4.1.1。
    \item \textbf{編譯器}：GCC / Clang。
    \item \textbf{關鍵資料結構庫}：使用 \texttt{stb\_ds.h} 實作高效能的動態陣列與雜湊表（Hash Map）。
\end{itemize}

\section{系統規格與邊界限制}

為確保組譯器在處理各種組合語言檔案時的穩定性，並防範潛在的緩衝區溢位（Buffer Overflow）等記憶體安全問題，本專案明確定義了系統的處理邊界，並制定了嚴謹的字元識別規範。

\subsection{系統邊界條件}

各項長度上限皆統一定義於 \lstinline{constants.h} 中，具體規格如下：

\begin{itemize}
    \item \textbf{單行讀取緩衝區（\lstinline{READ_BUFFER_SIZE}）}：
          配置 1024 Bytes，確保能安全讀取包含極長註解的單行原始碼。
    \item \textbf{標籤長度上限（\lstinline{MAX_LABEL_LENGTH}）}：
          64 字元。
    \item \textbf{助記符長度上限（\lstinline{MAX_MNEMONIC_LENGTH}）}：
          16 字元。
    \item \textbf{運算元長度上限（\lstinline{MAX_OPERAND_LENGTH}）}：
          256 字元，足以應付極長的字串常數宣告（如 \lstinline{C'...'}）。
    \item \textbf{T 紀錄長度上限（\lstinline{MAX_TRECORD_LENGTH}）}：
          30 Bytes。
          此數值嚴格遵循 SIC/XE 目的檔規範，確保單筆 Text Record 的機器碼長度不超過十六進位的 \lstinline{1E}。
\end{itemize}

\subsection{大小寫識別與處理}

在詞法分析與符號建構的過程中，系統針對不同型別的語彙基元（Token）採取以下判定策略：

\begin{itemize}
    \item \textbf{標籤與符號（Label \& Symbol）}：
          視為大小寫敏感（Case-sensitive）。
          例如，標籤 \lstinline{LOOP} 與 \lstinline{loop} 會被符號表（\lstinline{SYMTAB}）視為兩個完全不同的獨立參照。
    \item \textbf{助記符與假指令（Mnemonic \& Directives）}：
          要求嚴格匹配 \lstinline{OPTAB} 中定義的大寫字元，以維持組合語言指令結構的統一性。
\end{itemize}